% =========================================================================
% SECCIONES PRELIMINARES: ÍNDICES, ABREVIATURAS, SÍMBOLOS Y GLOSARIO
% =========================================================================

% 1. ÍNDICE GENERAL
\newpage
\phantomsection
\addcontentsline{toc}{section}{Índice general}
\tableofcontents 

% 2. ÍNDICE DE FIGURAS
\newpage
\phantomsection
\addcontentsline{toc}{section}{Índice de figuras}
\renewcommand{\listfigurename}{Índice de figuras}
\listoffigures

% 3. ÍNDICE DE TABLAS
\newpage
\phantomsection
\addcontentsline{toc}{section}{Índice de tablas}
\renewcommand{\listtablename}{Índice de tablas}
\renewcommand{\tablename}{Tabla}
\listoftables

% 4. ÍNDICE DE ALGORITMOS
\newpage
\phantomsection
\addcontentsline{toc}{section}{Índice de algoritmos}
\listof{algoritmo}{Índice de algoritmos}

% 5. LISTA DE ABREVIATURAS Y ACRÓNIMOS
\newpage
\phantomsection
\addcontentsline{toc}{section}{Lista de abreviaturas}
\section*{Lista de abreviaturas}

\begin{table}[H]
\centering
\small
\begin{tabularx}{\textwidth}{@{}p{3.2cm}X@{}}
\toprule
\rowcolor{tcAzul!12}
\textbf{Abreviatura / Sigla} & \textbf{Significado y Descripción Técnica} \\
\midrule
\textbf{AMLQ} & Academia Mayor de la Lengua Quechua (Cusco, Perú). \\
\textbf{BLEU} & \textit{Bilingual Evaluation Understudy} (métrica de concordancia de n-gramas léxicos). \\
\textbf{chrF++} & \textit{Character $n$-gram F-score with word $n$-grams} (métrica estándar para morfología aglutinante). \\
\textbf{CPU} & \textit{Central Processing Unit} (unidad central de procesamiento). \\
\textbf{DSR} & \textit{Design Science Research} (marco metodológico de investigación en ciencias del diseño). \\
\textbf{EIB} & Educación Intercultural Bilingüe. \\
\textbf{FAISS} & \textit{Facebook AI Similarity Search} (biblioteca de búsqueda vectorial densa eficiente). \\
\textbf{GGUF} & \textit{GPT-Generated Unified Format} (formato binario para cuantización e inferencia local eficiente). \\
\textbf{GPU} & \textit{Graphics Processing Unit} (unidad de procesamiento gráfico acelerado). \\
\textbf{HE} & Hipótesis Específica (HE1, HE2, HE3). \\
\textbf{HG} & Hipótesis General de la investigación. \\
\textbf{LLM} & \textit{Large Language Model} (modelo de lenguaje de gran escala, $>7$B parámetros). \\
\textbf{LoRA} & \textit{Low-Rank Adaptation} (técnica de adaptación mediante descomposición de bajo rango). \\
\textbf{NF4} & \textit{NormalFloat 4-bit} (tipo de dato no linealmente espaciado óptimo para cuantización de pesos). \\
\textbf{ODS} & Objetivos de Desarrollo Sostenible de las Naciones Unidas (ODS 9 y ODS 10). \\
\textbf{OE} & Objetivo Específico (OE1, OE2, OE3). \\
\textbf{OG} & Objetivo General de la investigación. \\
\textbf{PE} & Problema Específico (PE1, PE2, PE3). \\
\textbf{PEFT} & \textit{Parameter-Efficient Fine-Tuning} (ajuste fino eficiente en parámetros). \\
\textbf{PG} & Problema General de la investigación. \\
\textbf{PLN} & Procesamiento de Lenguaje Natural (equivalente a NLP en inglés). \\
\textbf{PPL} & Perplejidad (\textit{Perplexity}, métrica intrínseca de incertidumbre probabilística). \\
\textbf{PTQ} & \textit{Post-Training Quantization} (cuantización post-entrenamiento de pesos). \\
\textbf{QLoRA} & \textit{Quantized Low-Rank Adaptation} (adaptadores de bajo rango sobre pesos base en 4 bits). \\
\textbf{RAG} & \textit{Retrieval-Augmented Generation} (generación aumentada por recuperación). \\
\textbf{RAFT} & \textit{Retrieval-Augmented Fine-Tuning} (ajuste fino aumentado por recuperación). \\
\textbf{RAM} & \textit{Random Access Memory} (memoria principal del sistema). \\
\textbf{SLM} & \textit{Small Language Model} (modelo de lenguaje compacto, $\le 3$B parámetros). \\
\textbf{TPS} & \textit{Tokens Per Second} (rendimiento de generación de tokens por segundo). \\
\textbf{VRAM} & \textit{Video Random Access Memory} (memoria gráfica de video dedicada). \\
\textbf{WER} & \textit{Word Error Rate} (tasa de error de palabra en reconocimiento de texto/OCR). \\
\bottomrule
\end{tabularx}
\end{table}

% 6. LISTA DE SÍMBOLOS Y NOTACIÓN MATEMÁTICA
\newpage
\phantomsection
\addcontentsline{toc}{section}{Lista de símbolos}
\section*{Lista de símbolos}

\begin{table}[H]
\centering
\small
\begin{tabularx}{\textwidth}{@{}p{3.2cm}X@{}}
\toprule
\rowcolor{tcAzul!12}
\textbf{Símbolo / Notación} & \textbf{Definición y Rol Computacional} \\
\midrule
$W_0$ & Matriz de pesos preentrenada original del SLM congelada en 4 bits ($W_0 \in \mathbb{R}^{d \times k}$). \\
$\Delta W$ & Matriz de actualización adaptativa acumulada sobre los pesos del modelo ($\Delta W \in \mathbb{R}^{d \times k}$). \\
$A$ & Matriz de proyección de bajo rango inicializada de forma gaussiana ($A \in \mathbb{R}^{r \times k}$). \\
$B$ & Matriz de proyección de bajo rango inicializada en ceros ($B \in \mathbb{R}^{d \times r}$). \\
$r$ & Rango intrínseco de los adaptadores LoRA ($r \ll \min(d, k)$, fijado empíricamente en $r=16$). \\
$\alpha$ & Factor constante de escalamiento para la actualización LoRA ($\alpha = 32$). \\
$x$ & Vector de activación de entrada a la capa Transformer (tokens de consulta + contexto RAG). \\
$h$ & Vector de activación adaptado de salida: $h = W_0 x + \frac{\alpha}{r} (B \cdot A) x$. \\
$Q, K, V$ & Matrices de Consulta (\textit{Query}), Clave (\textit{Key}) y Valor (\textit{Value}) en las capas de autoatención. \\
$W_q, W_k, W_v, W_o$ & Matrices de proyección lineal de los cabezales de atención objetivo del ajuste fino QLoRA. \\
$d_k$ & Dimensión de escalamiento de los vectores de clave en el producto escalar de autoatención. \\
$\mathbf{e}_q$ & Vector de incrustación denso (\textit{embedding}) de la consulta formulada por el usuario. \\
$\mathbf{e}_d$ & Vector de incrustación denso del fragmento documental indexado en la base de conocimiento. \\
$\text{sim}(q, d)$ & Similitud del coseno normalizada entre los vectores de consulta y fragmento documental. \\
$\mathcal{D}$ & Colección curada de fragmentos documentales bilingües indexados en la base vectorial local. \\
$\mathcal{C}_k$ & Conjunto de los $k$ fragmentos documentales de mayor afinidad semántica recuperados. \\
$P_{\text{aug}}$ & Secuencia de tokens del prompt estructurado aumentado (instrucción + contexto + consulta). \\
$PPL$ & Perplejidad intrínseca calculada sobre el conjunto de evaluación en Quechua Collao. \\
\bottomrule
\end{tabularx}
\end{table}

% 7. ÍNDICE DE ECUACIONES
\newpage
\phantomsection
\addcontentsline{toc}{section}{Índice de ecuaciones}
\section*{Índice de ecuaciones}

\begin{table}[H]
\centering
\small
\begin{tabularx}{\textwidth}{@{}p{2.2cm}p{9.0cm}c@{}}
\toprule
\rowcolor{tcAzul!12}
\textbf{Ecuación} & \textbf{Descripción Matemática y Teórica} & \textbf{Página} \\
\midrule
Ecuación (1) & Descomposición de bajo rango para adaptación eficiente de pesos (LoRA / QLoRA) & \pageref{fig:esquema_qlora} \\
Ecuación (2) & Proyecciones lineales de Consulta ($Q$), Clave ($K$) y Valor ($V$) en capas Transformer & \pageref{fig:esquema_qlora} \\
Ecuación (3) & Mecanismo de autoatención con producto escalar escalado (\textit{Scaled Dot-Product Attention}) & \pageref{fig:esquema_qlora} \\
Ecuación (4) & Proyección lineal de salida multicabezal del bloque de atención Transformer & \pageref{fig:esquema_qlora} \\
Ecuación (5) & Métrica de similitud del coseno para recuperación densa de fragmentos documentales & \pageref{fig:pipeline_rag} \\
\bottomrule
\end{tabularx}
\end{table}

% 8. ÍNDICE DE ANEXOS
\newpage
\phantomsection
\addcontentsline{toc}{section}{Índice de anexos}
\section*{Índice de anexos}

\begin{table}[H]
\centering
\small
\begin{tabularx}{\textwidth}{@{}p{2.6cm}p{8.5cm}c@{}}
\toprule
\rowcolor{tcAzul!12}
\textbf{Identificador} & \textbf{Título y Contenido del Anexo} & \textbf{Página} \\
\midrule
Anexo A & Fichas de análisis y extracción crítica de antecedentes científicos (Fichas 1 a 6) & \pageref{sec:anexos} \\
Anexo B & Matriz de consistencia metodológica del proyecto de tesis (Problemas, Objetivos, Hipótesis) & \pageref{tab:matriz_consistencia} \\
Anexo C & Banco de 100 consultas patrón de prueba con contexto de referencia (\textit{ground-truth}) & \pageref{tab:metricas} \\
\bottomrule
\end{tabularx}
\end{table}

% 9. GLOSARIO DE TÉRMINOS
\newpage
\phantomsection
\addcontentsline{toc}{section}{Glosario de términos}
\section*{Glosario de términos}

\begin{description}[leftmargin=1.8em, itemsep=0.6em]
	\item[\textbf{Aglutinación:}] Propiedad morfológica del Quechua Collao donde las palabras se componen de una raíz léxica invariable a la cual se concatenan múltiples sufijos ordenados jerárquicamente, condensando funciones sintácticas extensas en vocablos individuales.
	\item[\textbf{Alucinación semántica:}] Fenómeno en el cual un modelo de lenguaje generativo produce afirmaciones gramaticalmente fluidas pero fácticamente erróneas, no sustentadas en los datos de entrada o contrarias a la realidad empírica.
	\item[\textbf{Cuantización (GGUF / NF4):}] Técnica de compresión numérica post-entrenamiento que reduce la precisión de los pesos de la red neuronal (de punto flotante de 16 bits a enteros o cuantizaciones no lineales de 4 bits), disminuyendo drásticamente el uso de memoria RAM/VRAM con mínima pérdida de perplejidad.
	\item[\textbf{Design Science Research (DSR):}] Marco metodológico riguroso de ingeniería en informática y sistemas orientado a la creación iterativa y evaluación formal de artefactos tecnológicos diseñados para solucionar problemas prácticos específicos.
	\item[\textbf{Dispersión atencional (\textit{Attention Dispersion}):}] Incapacidad observada en modelos compactos no adaptados para concentrar la masa de probabilidad de autoatención sobre los fragmentos documentales relevantes inyectados por el RAG, provocando la omisión del contexto.
	\item[\textbf{Fidelidad factual (\textit{Faithfulness}):}] Métrica computacional objetiva (del marco Ragas) que cuantifica la proporción de afirmaciones generadas por el modelo que pueden derivarse e inferirse lógicamente a partir del contexto recuperado.
	\item[\textbf{Generación Aumentada por Recuperación (RAG):}] Arquitectura computacional que complementa la generación de un modelo de lenguaje vinculándolo a un almacén documental externo dinámico mediante búsqueda vectorial por similitud semántica.
	\item[\textbf{Inferencia local \textit{on-device}:}] Ejecución autónoma y desatendida de un modelo de aprendizaje profundo directamente en el hardware del usuario final (CPU o GPU modesta) sin requerir transmisión de datos ni llamadas a APIs en la nube.
	\item[\textbf{NormalFloat4 (NF4):}] Tipo de cuantización no lineal estadísticamente óptimo para pesos de redes neuronales distribuidos normalmente con media cero, base fundamental del método QLoRA.
	\item[\textbf{QLoRA (\textit{Quantized Low-Rank Adaptation}):}] Método de ajuste fino eficiente que congela el modelo base cuantizado en 4 bits e inyecta pares de matrices entrenables de bajo rango ($A$ y $B$) en las capas de atención, reduciendo en más del 80\% el consumo de VRAM.
	\item[\textbf{Quechua Collao:}] Variante dialectal de la familia Quechua Sureño hablada en las regiones de Cusco y Puno, caracterizada por un inventario fonológico con oclusivas simples, glotizadas y aspiradas, y una ortografía oficial normalizada por el Ministerio de Educación del Perú.
	\item[\textbf{Small Language Model (SLM):}] Modelo de lenguaje compacto con una escala paramétrica inferior o igual a 3 mil millones de parámetros ($\le 3$B), diseñado para ofrecer inferencia interactiva sobre computadoras personales o dispositivos de borde.
\end{description}